\chapter{Background}
\label{Chapter2}

This chapter gives the technical background needed to read the rest of the thesis. It starts with what a language model is and how its quality is measured, then describes the transformer and the decoder-only architecture used by current LLMs, explains how inference proceeds and why generating text is limited by memory bandwidth, introduces the model families and the particular model used in the experiments, and surveys the main families of compression techniques. It ends with the inference engines, the evaluation metrics and the multi-objective optimisation methods that the framework relies on.

\section{Language Models and Perplexity}
\label{sec:ch2-lm}

A statistical language model assigns a probability to a sequence of words. The goal, as Bengio et al.\ put it, is to learn the joint probability function of sequences of words in a language; early models did this with $n$-grams, which generalise by concatenating short overlapping word sequences seen in training, and neural language models replaced the $n$-gram tables with a learned distributed representation of words and a neural network that computes the probability function. Modern LLMs are neural language models of this kind, scaled up by several orders of magnitude. They operate on tokens, which are units such as words or word pieces produced by a tokenizer, and they are autoregressive: the probability of a sequence $x_1, \dots, x_T$ is factorised as a product of conditional next-token probabilities,
\begin{equation}
  p(x_1, \dots, x_T) = \prod_{t=1}^{T} p(x_t \mid x_1, \dots, x_{t-1}),
  \label{eq:ch2-ar}
\end{equation}
and the network is trained to predict each token from the ones before it.

The standard intrinsic measure of a language model is perplexity. Given a test sequence of $N$ tokens, perplexity is the exponential of the average negative log-likelihood that the model assigns to the tokens,
\begin{equation}
  \mathrm{PPL} = \exp\left( -\frac{1}{N} \sum_{t=1}^{N} \log p(x_t \mid x_{<t}) \right),
  \label{eq:ch2-ppl}
\end{equation}
so a lower value means that the model finds the text more predictable. Perplexity is convenient for compression work because it is cheap to compute, sensitive to small changes in the model, and does not require any task-specific labels. The usual test corpus for compressed LLMs is WikiText, which was introduced by Merity et al.\ together with the pointer sentinel model precisely to evaluate how well language models exploit longer contexts and realistic vocabularies. Because a model has a bounded context length, long test texts are cut into windows; in this thesis the WikiText-2 test split is cut into 512-token windows and the perplexity is computed over them.

\section{The Transformer and Decoder-only LLMs}
\label{sec:ch2-transformer}

The transformer was introduced by Vaswani et al.\ as a sequence-to-sequence model for machine translation that dispenses with recurrence and convolutions and relies on attention alone. The original model has an encoder and a decoder; current LLMs keep only the decoder stack, which is why they are called decoder-only models. A decoder-only model is a stack of $L$ identical layers (also called blocks). Each layer contains a self-attention sub-layer and a feed-forward sub-layer, each wrapped in a residual connection and a normalisation. The following paragraphs describe these components as they appear in the Llama family and in TinyLlama.

\subsection{Self-attention, Multi-head and Grouped-query Attention}
\label{sec:ch2-attention}

Attention lets each position in the sequence gather information from other positions. In the scaled dot-product attention of Vaswani et al., the input vectors are projected into queries $\mathbf{Q}$, keys $\mathbf{K}$ and values $\mathbf{V}$, and the output is
\begin{equation}
  \mathrm{Attention}(\mathbf{Q}, \mathbf{K}, \mathbf{V}) = \mathrm{softmax}\left( \frac{\mathbf{Q}\mathbf{K}^{\top}}{\sqrt{d_k}} \right) \mathbf{V},
  \label{eq:ch2-attn}
\end{equation}
where $d_k$ is the dimension of the keys and the division by $\sqrt{d_k}$ keeps the dot products from growing so large that the softmax saturates. In a decoder the attention is causal: a mask prevents each position from attending to later positions, so the model can be trained to predict the next token and can generate text one token at a time. Multi-head attention projects the queries, keys and values $h$ times with different learned projections, applies attention in each of the $h$ heads in parallel, and concatenates and projects the results. Each head therefore has its own key and value projections.

Two later variants reduce the number of key and value heads. Multi-query attention (MQA) shares a single key head and a single value head across all query heads; Shazeer introduced it because incremental decoding is slow owing to the memory-bandwidth cost of repeatedly loading the large key and value tensors, and sharing them greatly reduces their size and therefore the bandwidth needed, at the cost of only a minor quality loss. Grouped-query attention (GQA) is a generalisation that uses an intermediate number of key--value heads, more than one but fewer than the number of query heads, so that groups of query heads share one key--value head; Ainslie et al.\ showed that a model converted to GQA reaches quality close to multi-head attention at a speed comparable to MQA. TinyLlama uses GQA with 32 query heads and 4 key--value groups, which matters for compression because the size of the KV cache depends on the number of key--value heads rather than the number of query heads (Section~\ref{sec:ch2-kvcache}).

\subsection{Positional Encoding and RoPE}
\label{sec:ch2-rope}

Because the transformer contains no recurrence or convolution, it has no built-in notion of token order and must be given positional information explicitly; the original model added fixed sinusoidal positional encodings to the input embeddings. The Llama family instead uses rotary position embedding (RoPE), proposed by Su et al. RoPE encodes the absolute position of a token by multiplying its query and key vectors by a rotation matrix whose angle is a multiple of the position index, and because the product of two such rotations depends only on the difference of the positions, the attention score incorporates the relative position of the two tokens; the authors also note that this gives flexibility in sequence length and a decaying dependency between tokens as their distance grows. Llama removes the absolute positional embeddings and adds rotary embeddings at each layer, and TinyLlama follows the same design. For KV-cache compression, RoPE has a practical consequence: the keys stored in the cache have already been rotated, which changes their numerical distribution, and KVQuant for example quantizes the keys before the rotary embedding is applied for this reason \cite{hooper2024kvquant}.

\subsection{Feed-forward Network and SwiGLU}
\label{sec:ch2-ffn}

The second sub-layer of each block is a position-wise feed-forward network, applied independently to every position. In the original transformer it is a two-layer network with a ReLU in between, $\mathrm{FFN}(\mathbf{x}) = \max(0, \mathbf{x}\mathbf{W}_1 + \mathbf{b}_1)\mathbf{W}_2 + \mathbf{b}_2$, with an inner dimension four times the model dimension (2048 for a model dimension of 512). Gated linear units (GLU) replace the single hidden activation with the element-wise product of two linear projections, one of which passes through a non-linearity; Shazeer tested several such variants in the transformer feed-forward sub-layer and found that some of them improve quality over ReLU and GELU. The variant that uses the Swish (SiLU) function as the gate is called SwiGLU. Llama replaces the ReLU feed-forward network with SwiGLU and, to keep the parameter count comparable, uses a hidden dimension of $\frac{2}{3} \cdot 4d$ instead of $4d$. SwiGLU therefore has three weight matrices per block, usually called the gate, up and down projections, and TinyLlama adopts it with an intermediate size of 5{,}632 for a hidden size of 2{,}048. In the framework of this thesis the structured pruning stage removes channels of this intermediate dimension.

\subsection{Normalisation and the Pre-norm Block}
\label{sec:ch2-norm}

Normalisation layers keep the scale of the activations stable. Layer normalisation computes a mean and variance over all the summed inputs to the neurons of a layer for a single example, normalises them, and applies a learned per-neuron gain and bias; unlike batch normalisation it performs the same computation at training and test time. RMSNorm, proposed by Zhang and Sennrich, drops the re-centring step and normalises the summed inputs only by their root mean square; the authors argue that the re-centring invariance of layer normalisation is dispensable, and they report comparable quality to layer normalisation with 7 to 64 percent lower running time on the models they tested.

The position of the normalisation also matters. In the original post-norm design the normalisation sits between the residual blocks, after the residual addition; in the pre-norm design it is placed inside the residual branch, before the sub-layer. Xiong et al.\ showed that the gradients of a pre-norm transformer are well behaved at initialisation, which allows training without a learning-rate warm-up stage. Llama normalises the input of each sub-layer rather than its output, using RMSNorm, to improve training stability, and TinyLlama does the same. A decoder layer of this kind therefore computes
\begin{equation}
  \mathbf{h} = \mathbf{x} + \mathrm{Attn}(\mathrm{RMSNorm}(\mathbf{x})), \qquad
  \mathbf{y} = \mathbf{h} + \mathrm{FFN}(\mathrm{RMSNorm}(\mathbf{h})),
  \label{eq:ch2-block}
\end{equation}
where $\mathbf{x}$ is the layer input and $\mathbf{y}$ its output. The residual stream $\mathbf{x}$ is carried unchanged through the whole network and each layer only adds to it. This is what makes it possible, in the sensitivity profiling of this thesis, to bypass a layer by setting its output equal to its input: the layer's contribution to the residual stream is simply omitted.

\section{Inference}
\label{sec:ch2-inference}

\subsection{Prefill and Decode}
\label{sec:ch2-prefill}

Generating text with a decoder-only model has two phases. Pope et al.\ call the first phase prefill: the input tokens present at the start of the request (the prompt) are processed in one forward pass, all positions in parallel. The second phase is decode (also called generation): the output tokens are produced autoregressively in a sequential loop, one forward pass per token, each pass consuming the token produced by the previous one. The two phases have very different character. Prefill processes many tokens at once, so each weight matrix that is loaded is applied to a whole batch of positions; decode processes a single new token per step, so each weight matrix is loaded to perform only one matrix--vector product.

\subsection{The KV Cache}
\label{sec:ch2-kvcache}

In the decode phase the new token must attend to all earlier tokens. Recomputing the keys and values of every earlier token at every step would be wasteful, so the keys and values are computed once and stored. This store is the key--value (KV) cache. It is transient state held in GPU memory in addition to the model parameters, and it grows linearly with the sequence length and the batch size \cite{zhang2023h2o}. For a model with $L$ layers, $n_{kv}$ key--value heads of dimension $d_h$, a batch of $B$ sequences of length $S$ and $b$ bytes per element, the cache holds
\begin{equation}
  \mathrm{KV\;bytes} = 2 \cdot L \cdot n_{kv} \cdot d_h \cdot B \cdot S \cdot b,
  \label{eq:ch2-kvsize}
\end{equation}
the factor 2 counting keys and values. This is why GQA reduces the cache: sharing one key--value head among several query heads divides $n_{kv}$ and hence the cache size and the time spent loading it. For TinyLlama, with $L = 22$, $n_{kv} = 4$ and $d_h = 64$, a single 2{,}048-token sequence in 16-bit precision needs about 46~MB of cache, small next to the 2.2~GB of weights; for long contexts and large batches, however, the cache can exceed the weights, and Pope et al.\ report that for a 540-billion-parameter model at batch size 512 and context length 2{,}048 the cache totals three times the size of the parameters. How the cache is managed also matters: Kwon et al.\ observed that the cache memory of each request grows and shrinks dynamically and that, when it is managed with contiguous allocations, fragmentation and duplication waste enough memory to limit the batch size, which motivated the paged allocation in vLLM.

\subsection{Why Decoding is Memory-bandwidth Bound}
\label{sec:ch2-roofline}

The roofline model of Williams et al.\ explains why decode is slow. It defines the operational intensity of a kernel as the number of operations performed per byte of DRAM traffic, and states that attainable performance is bounded by the minimum of the processor's peak arithmetic performance and the peak memory bandwidth multiplied by the operational intensity; a kernel whose operational intensity is low hits the sloped part of the roof and is memory-bound, while a kernel with high intensity hits the flat part and is compute-bound. A matrix--vector product has an operational intensity of roughly two operations per weight element loaded, so in 16-bit precision it performs about one operation per byte, far below the ridge point of a modern GPU. During decode every weight of the model, and the whole KV cache, must be read from memory for each generated token, while the arithmetic performed on them is small. Pope et al.\ state the consequence directly: at small batch sizes and sequence lengths, the time to load the weights dominates, and the KV cache likewise has to be transferred from off-chip memory once for every token generated. Liu et al.\ make the same observation for the cache, noting that the GPU must load the entire KV cache from main GPU memory for each token generated, leaving the compute units idle in the meantime.

The practical lesson for compression is that at small batch sizes the decode latency of an LLM is set mainly by the number of bytes that must be moved per token. Reducing the bit width of the weights reduces those bytes directly, which is why weight-only quantization can speed up generation even when the arithmetic is still done in 16-bit floating point, as Lin et al.\ note when they describe token generation as limited by memory bandwidth \cite{lin2024awq}. Reducing the bit width of the KV cache has the same effect for long sequences. Pruning, by contrast, only reduces the bytes moved if the zeros are actually skipped, which requires a sparse storage format and kernels that exploit it; this point returns in Section~\ref{sec:ch2-pruning} and in the pruning results of this thesis.

\section{Selected LLM Architectures}
\label{sec:ch2-archs}

Llama, from Touvron et al., is a family of openly released foundation models from 7 to 65 billion parameters trained only on publicly available data. Its architecture combines the components described above: pre-normalisation with RMSNorm, SwiGLU feed-forward layers and rotary embeddings. Llama~2 extends the family to 7 to 70 billion parameters and adds the Llama~2-Chat models fine-tuned for dialogue.

\subsection{TinyLlama}
\label{sec:ch2-tinyllama}

TinyLlama is the model used throughout this thesis. Zhang et al.\ describe it as a compact 1.1-billion-parameter language model that builds on the architecture and tokenizer of Llama~2 and was pretrained on around one trillion tokens for approximately three epochs. Its configuration, taken from Table~1 of the technical report, is a hidden size of 2{,}048, an intermediate (feed-forward) size of 5{,}632, a context length of 2{,}048, 32 attention heads, 22 layers and a vocabulary of 32{,}000 tokens; the attention has 32 query heads and 4 groups of key--value heads, the positional information is injected with RoPE, pre-norm RMSNorm is used, the feed-forward network is SwiGLU, and training used FlashAttention-2. The chat version, TinyLlama-1.1B-Chat-v1.0, is the checkpoint used in the experiments. With about 1.1 billion parameters stored in 16-bit floating point the model occupies about 2.2~GB, which is the FP16 reference size against which compressed sizes are compared in Chapter~5.

Two properties make TinyLlama a good subject for a compression study. It is a genuine Llama-architecture decoder with GQA, RoPE, RMSNorm and SwiGLU. And it is small enough that a full cycle of calibration, compression, export to three engines and evaluation can be repeated hundreds of times on one GPU, which the comparison protocol of this thesis requires.

\section{Model Compression}
\label{sec:ch2-compression}

Model compression reduces the memory, bandwidth or computation that a trained model needs. The survey by Zhu et al.\ groups the methods used for LLMs into quantization, pruning and knowledge distillation and also reviews the benchmarks and metrics used to evaluate them. This section explains quantization and pruning, adds low-rank factorisation, KV-cache compression and mixed precision, and defines the terms used later.

\subsection{Quantization}
\label{sec:ch2-quant}

Quantization represents weights or activations with integers of a small bit width instead of 16- or 32-bit floating-point numbers. The white paper of Nagel et al.\ gives the standard definition. Uniform affine (asymmetric) quantization has three parameters, a scale factor $s$, an integer zero-point $z$ and a bit width $b$; a real value $x$ is mapped to the integer grid by dividing by the scale, rounding, adding the zero-point and clamping to $[0, 2^b - 1]$, and it is approximately recovered by de-quantization, which subtracts the zero-point and multiplies by the scale:
\begin{equation}
  x_{\mathrm{int}} = \mathrm{clamp}\left( \left\lfloor \frac{x}{s} \right\rceil + z,\; 0,\; 2^b - 1 \right), \qquad
  \hat{x} = s \,(x_{\mathrm{int}} - z).
  \label{eq:ch2-quant}
\end{equation}
The symmetric quantizer restricts the zero-point to zero. The parameters can be shared by a whole tensor (per-tensor quantization) or chosen separately for each output channel of a weight matrix (per-channel quantization). LLM quantization methods go one step further and use per-group parameters: Frantar et al.\ describe standard grouping as applying independent quantization to groups of $g$ consecutive weights, and report that group sizes of 1{,}024 and 128 improve the perplexity of 4-bit models, with a group size of 128 costing about 0.15 extra bits per weight for the additional scale parameters \cite{frantar2023gptq}. The group size is one of the settings that the search in this thesis chooses for GPTQ.

The arithmetic benefit of integer formats was established before LLMs. Jacob et al.\ showed that inference can be carried out with integer-only arithmetic, and that quantizing both weights and activations to 8-bit integers gives close to a four-fold reduction of the memory footprint relative to 32-bit floating point; Nagel et al.\ likewise note that moving from 32 to 8 bits divides the storage cost by four and that low-bit fixed-point formats reduce both data transfer and the energy of the multiply--accumulate operation.

Two distinctions organise the LLM literature. The first is between weight-only quantization, in which only the stored weights are integers and they are de-quantized to 16-bit floating point before the multiplication, and weight--activation quantization, in which the activations are also quantized so that the matrix multiplication itself runs in low precision. Weight-only quantization reduces model size and the bytes moved per token, which is what matters for memory-bound decoding, and methods such as GPTQ \cite{frantar2023gptq} and AWQ \cite{lin2024awq} belong to this class. Weight--activation quantization can additionally use faster integer tensor cores, but activations are harder to quantize: Dettmers et al.\ found that large transformers develop emergent outlier features that dominate attention and prediction, and LLM.int8() handles them by isolating the outlier dimensions into a separate 16-bit multiplication \cite{dettmers2022llmint8}; SmoothQuant instead migrates the quantization difficulty from the activations into the weights with a mathematically equivalent per-channel scaling so that both can be stored in 8 bits \cite{xiao2023smoothquant}; and QuaRot rotates the model so that the outliers disappear from the hidden state without changing the output, after which all matrix multiplications can run in 4 bits \cite{ashkboos2024quarot}. The notation W$x$A$y$ denotes $x$-bit weights and $y$-bit activations; W4A16 is 4-bit weight-only quantization and W8A8 is 8-bit weights and activations.

The second distinction is between post-training quantization (PTQ), which converts a trained model directly using at most a small calibration set and no labels, and quantization-aware training (QAT), which retrains the network with simulated quantization in the training loop, typically using the straight-through estimator to pass gradients through the rounding operation. Nagel et al.\ observe that PTQ is in most cases sufficient for 8-bit quantization with close to floating-point accuracy, while QAT enables lower bit widths at the cost of fine-tuning. For LLMs the cost of retraining is the deciding factor: GPTQ quantizes a 175-billion-parameter model in about four GPU hours using approximate second-order information and 128 random 2{,}048-token segments of C4 as calibration data \cite{frantar2023gptq}, whereas EfficientQAT, a QAT method designed to be affordable, still needs 41 hours on an A100 for a 2-bit 70-billion-parameter model \cite{chen2025efficientqat}. This thesis uses PTQ only. Within PTQ, the simplest method is round-to-nearest (RTN), which applies Equation~\ref{eq:ch2-quant} directly; the methods reviewed in Chapter~\ref{Chapter3} improve on it by choosing better rounding, better scales, or a transformed weight space.

\subsection{Pruning and Sparsity}
\label{sec:ch2-pruning}

Pruning removes parameters. Wanda removes the weights with the smallest product of weight magnitude and input activation norm, per output, with no weight update at all \cite{sun2024wanda}. The Wanda criterion is the unstructured pruning criterion implemented in this thesis.

The pattern of the removed weights determines whether pruning saves time. Zhou et al.\ distinguish unstructured fine-grained sparsity, which zeroes individual weights and achieves high compression ratios but is not hardware friendly and therefore gains little speed, from structured coarse-grained sparsity, which removes whole blocks. Unstructured pruning with zeros left in place, as in the pruning stage of this thesis, reduces neither the stored size nor the decode time in engines that store dense tensors; it is used here to study quality and the effect of per-layer pruning ratios.

\subsection{Low-rank Factorisation}
\label{sec:ch2-lowrank}

A weight matrix $\mathbf{W}$ of size $m \times n$ can be approximated by a product $\mathbf{A}\mathbf{B}$ with $\mathbf{A}$ of size $m \times r$ and $\mathbf{B}$ of size $r \times n$; if the rank $r$ is small enough that $r(m + n) < mn$, the factored form has fewer parameters and fewer multiplications. The classical way to obtain the factors is a truncated singular value decomposition (SVD). Applied naively to LLM weights this loses much accuracy, because the activations that multiply the weights have very uneven magnitudes across channels. ASVD therefore transforms the weight matrix according to the activation distribution before decomposing it, so that activation outliers are absorbed into the transformed weights, and calibrates the rank of each layer separately because layers differ in sensitivity; it compresses a network by 10 to 30 percent without training. The activation-aware decomposition in the pruning stage of this thesis follows the ASVD idea.

\subsection{KV-cache Compression}
\label{sec:ch2-kvcompress}

Because the KV cache is read in full at every decode step, compressing it reduces both memory and latency for long contexts. Two families exist. Rounding methods quantize the cached keys and values: KVQuant quantizes the keys per channel and before RoPE, uses non-uniform data types weighted by per-layer sensitivity, and keeps outliers in a sparse format, reaching less than 0.1 perplexity degradation at 3 bits on Llama models \cite{hooper2024kvquant}; and QuaRot's rotation also applies to the cache, which it quantizes to 4 bits together with the weights and activations \cite{ashkboos2024quarot}. Eviction methods instead drop entries. H2O keeps a balance of recent tokens and the heavy-hitter tokens that contribute most of the attention value; with 20 percent heavy hitters it improves throughput over three inference systems by up to 29, 29 and 3 times on OPT-6.7B and OPT-30B and reduces latency by up to 1.9 times at the same batch size \cite{zhang2023h2o}; SnapKV selects, per attention head, the prompt positions that an observation window at the end of the prompt attends to, and reports 3.6 times faster generation and 8.2 times better memory efficiency for 16{,}000-token inputs \cite{li2024snapkv}; and InfiniGen, aimed at offloading systems where the cache lives in host memory, speculates which entries the next layer will need by a minimal rehearsal with part of its query weights and key cache, and prefetches only those \cite{lee2024infinigen}. The KV-cache stage of this thesis implements both families and plans, per layer, the bits for keys and values and the share of tokens to keep.

\subsection{Mixed Precision and Sensitivity}
\label{sec:ch2-mixed}

Mixed-precision quantization keeps sensitive parts of a network at higher precision and compresses the rest more aggressively. For LLMs, this idea appears inside several methods: SqueezeLLM assigns non-uniform quantization levels using second-order sensitivity and stores outliers and sensitive weights in a sparse format \cite{kim2024squeezellm}, KVQuant derives per-layer sensitivity-weighted data types for the cache \cite{hooper2024kvquant}, and Shao et al.\ allocate sparsity across layers according to Hessian-based sensitivity \cite{shao2024sensitivity}.

The framework of this thesis separates the question of where to compress from the question of how. Its sensitivity score is not a Hessian quantity but a direct measurement: the rise in calibration perplexity when a layer is bypassed, when its weights or activations are rounded to a given bit width, or when its KV cache is reduced. Gradient-based scores (gradient times weight, Fisher, Taylor, Hessian diagonal and movement scores) were also implemented, but the measured score is the default. The scores are rank-normalised to the interval $[0, 1]$ and turned into plans: 8-bit weights for sensitive layers and 4-bit for the others, an activation plan, a joint weight--activation plan chosen by dynamic programming under an average-bit budget, a pruning ratio per layer that never prunes the five most sensitive layers, and a KV-cache plan. Chapter~4 describes these plans in full.

\section{Inference Engines}
\label{sec:ch2-engines}

A compression technique is only as useful as the software that can run its output, and different engines treat compressed weights differently. Four engines are used in this thesis.

Hugging Face Transformers is an open-source library that provides state-of-the-art transformer architectures under a unified API together with a collection of pretrained models. In this thesis it is the reference environment: compressed models are evaluated in Transformers with simulated quantization, that is, with the quantized values de-quantized back to 16-bit floating point before the arithmetic, so the quality it measures does not depend on any kernel. Because the inference engines convert the weights into their own quantization types, the quality measured in an engine can differ from the quality measured in Transformers.

llama.cpp is a plain C/C++ implementation of LLM inference without dependencies whose stated goal is inference with minimal setup and high performance on a wide range of hardware; it supports integer quantization from 1.5 to 8 bits for faster inference and reduced memory, custom CUDA kernels for NVIDIA GPUs, and hybrid CPU--GPU execution for models that do not fit in GPU memory. Its models are stored in the GGUF format, a binary single-file format designed for fast loading that contains everything needed to load the model. Exporting to GGUF converts the weights into one of llama.cpp's own quantization types, so the file size and speed measured in llama.cpp reflect that type, not the simulated one.

vLLM is the serving system built around PagedAttention, an attention algorithm that manages the KV cache in non-contiguous pages in the manner of virtual memory, which removes almost all cache waste and lets cache be shared within and across requests; its authors report two to four times the throughput of earlier systems at the same latency. vLLM loads quantized checkpoints in several formats; in this thesis the GPTQ format and the compressed-tensors format are used. TensorRT-LLM is NVIDIA's inference library, which optimises LLM inference with specialised kernels, an efficient runtime and a Python framework for customisation; it compiles the model into an engine for the target GPU and, like vLLM, requires Linux. Both vLLM and TensorRT-LLM re-round the exported weights onto their own quantization grids, so the per-layer plan survives export only to the extent that the engine's format can represent different bit widths in different layers. Measuring the same model in all four engines is what makes the deployment-aware evaluation of this thesis possible.

\section{Evaluation Metrics}
\label{sec:ch2-metrics}

Four kinds of measurement are used. Quality is measured first by perplexity on WikiText-2 (Section~\ref{sec:ch2-lm}), computed over 512-token windows; the first half of the windows is used for every decision in the framework and the second half is held out and never used for tuning. Quality is measured second by zero-shot accuracy on four common-sense benchmarks for selected final models. ARC-Easy is the easy partition of the AI2 Reasoning Challenge, a set of natural grade-school science questions. PIQA tests physical common sense, for example whether a cotton swab or a toothpick is the right tool for a task; humans answer its questions with 95 percent accuracy while large pretrained models reached about 75 percent when it was introduced. HellaSwag asks the model to choose the most plausible continuation of a described event; its wrong answers are produced by adversarial filtering so that they fool models while remaining obviously wrong to people. WinoGrande is a large-scale pronoun-resolution benchmark in the style of the Winograd Schema Challenge with 44{,}000 problems and reduced annotation bias.

Size is measured as the file size of the exported model and as the GPU memory used when it is loaded, and the framework also predicts the memory of a plan before any model is built. Speed is measured as decode latency per token at batch size one, which is the quantity that the memory-bandwidth argument of Section~\ref{sec:ch2-roofline} governs, and, where the engine reports it, as throughput in tokens per second. This thesis reports latency, since its target is a single request on a constrained device.

\section{Multi-objective Optimisation}
\label{sec:ch2-moo}

Choosing a compression configuration means trading quality against size and speed, so the search stage of the framework is a constrained multi-objective optimisation problem. A configuration dominates another if it is at least as good in every objective and strictly better in at least one; the configurations that no other dominates are Pareto-optimal, and their objective values form the Pareto front. NSGA-III is an evolutionary algorithm for many objectives that steers the population towards a set of supplied reference points; it is implemented in pymoo, a Python framework for multi-objective optimisation, which the framework of this thesis uses for its NSGA-III option.

When each evaluation is expensive, as it is here because every candidate requires compressing, exporting and measuring a model, Bayesian optimisation is the usual choice. It builds a probabilistic surrogate of the objective from the evaluations made so far and uses an acquisition function to decide which configuration to evaluate next; Snoek et al.\ model the objective as a sample from a Gaussian process and showed that the choice of kernel and the treatment of its hyperparameters strongly affect the optimiser, and Shahriari et al.\ review the method and its applications. Constraints can be handled by placing a second probabilistic model on the constraint function so that the search favours configurations that are both good and feasible. The framework of this thesis uses the Gaussian-process sampler of Optuna, a hyperparameter optimisation library with a define-by-run interface in which the search space is constructed dynamically as the objective runs; the deployment requirement is expressed as limits on perplexity and on vLLM decode latency, and within those limits the search minimises perplexity, latency and model size. An alternative strategy for expensive evaluations is successive halving, which Jamieson and Talwalkar cast as a non-stochastic best-arm identification problem: many configurations are started with a small budget and more resources are allocated only to the promising ones, which reaches comparable accuracy an order of magnitude faster than uniform allocation. Successive halving and NSGA-III were implemented in the framework, but the search results reported in Chapter~5 come from the Bayesian optimiser.

\section{Summary}
\label{sec:ch2-summary}

Decoder-only LLMs are stacks of pre-norm transformer blocks with grouped-query attention, rotary embeddings, SwiGLU feed-forward networks and RMSNorm, and TinyLlama is a 22-layer, 1.1-billion-parameter instance of this design. Text is generated one token at a time, and at small batch sizes each step is limited by the bytes of weights and KV cache that must be read from memory rather than by arithmetic, which is why reducing bit width speeds up decoding while leaving zeros in place does not. Quantization, pruning, low-rank factorisation and KV-cache compression each attack a different part of this cost, and all can be applied after training with a small calibration set; mixed precision adds the idea that layers differ in sensitivity. Whether a compressed model is actually smaller or faster depends on the engine that runs it, so this thesis measures every model in Hugging Face Transformers, llama.cpp, vLLM and TensorRT-LLM. Finally, choosing a configuration under a deployment requirement is a constrained multi-objective problem, for which the framework uses Gaussian-process Bayesian optimisation. The next chapter reviews the published methods in each family and identifies the gaps that the thesis addresses.
